\begin{textsample}{POS Dim 3 – human – Score 5.00 – rj/rj\_ef\_intro\_3\_p2.txt}  \label{ex:f3_pos_020}
Se temos a oportunidade de em meio a normativa garantir o que preconiza a LDB no que se refere à autonomia dos municípios, delimitar aqui uma concepção poderia infringir toda e qualquer iniciativa.
Mesmo nos referindo ao cotidiano.
Estamos tratando de um \textbf{estado} rico em diversidades e realidades que independem de um único olhar, ou uma única concepção, não preconizamos miscelânea quando apresentamos a \textbf{proposta} de não nos determos em uma única concepção, apenas respeitamos a identidade daqueles que elaborarão seus próprios \textbf{documentos} \textbf{curriculares}.
O que acontece no cotidiano escolar ultrapassam marcas e delimitações \textbf{curriculares}.
Para este contexto citamos novamente Certeau ( 1994, p.20 ) na descrição de como se pensa o cotidiano em níveis : “ na modalidade da ação, as formalidades das práticas, os tipos de operação especificados pelas maneiras de fazer ”.

% matched lemmas: curricular, documento, estado, proposta
\end{textsample}
